\chapter{Единый квадрат кривизны M300 и гессиан связующего сектора}

Кинематическая склейка оставила открытым происхождение действия и
устойчивость двух массивных связующих комбинаций. Модулярная высота даёт
естественное уравнение плоской градуированной структуры
\begin{equation}
 [d,d^\dagger]=h.
\end{equation}
Поэтому проверим норму отклонения от этого уравнения:
\begin{equation}
 \mathcal F_H=[d,d^\dagger]-h,
 \qquad
 S_H=\tau_{300}(\mathcal F_H^2).
 \label{eq:v5-m300-hodge-action}
\end{equation}

\section{Три блока одного функционала}

На частичной трёхвершинной цепи
\begin{equation}
 d=E_{10}\otimes X+E_{21}\otimes\Phi I_3,
 \qquad h=\operatorname{diag}(-I_3,0,I_3),
\end{equation}
кривизна имеет блоки
\begin{equation}
 \mathcal F_H=\operatorname{diag}\left(
 I_3-X^\dagger X,
 XX^\dagger-|\Phi|^2I_3,
 (|\Phi|^2-1)I_3
 \right).
\end{equation}
После нормировки следа и сокращения одинакового множителя SU(5) получаем
\begin{equation}
 \begin{split}
 V_H(X,\Phi)=\frac13\{&
 \Tr(I_3-X^\dagger X)^2
 +\Tr(XX^\dagger-|\Phi|^2I_3)^2\\
 &+3(|\Phi|^2-1)^2\}.
 \end{split}
 \label{eq:v5-m300-reduced-potential}
\end{equation}
Таким образом, один квадрат одновременно содержит левое закрепление
семейного кадра, среднее отображение момента и правую нормировку
конденсата. Отдельные портальные коэффициенты не вводятся. Прямая проверка
матричной и сокращённой форм даёт максимальный остаток
\(1.5\cdot10^{-14}\).

\section{Точный гессиан при нулевом импульсе}

Нулевой минимум равен
\begin{equation}
 X_\star=I_3,
 \qquad |\Phi_\star|=1,
 \qquad V_H=0.
\end{equation}
На девяти вещественных компонентах $X$ и двух компонентах $\Phi$
точный спектр гессиана равен
\begin{equation}
 \Spec\Hess V_H=
 \left\{
 0^{(4)},
 \left(\frac{16}{3}\right)^{(5)},
 \frac{32-8\sqrt7}{3},
 \frac{32+8\sqrt7}{3}
 \right\}.
 \label{eq:v5-m300-zero-hessian}
\end{equation}
Четыре нулевых направления состоят из трёх касательных к орбите SO(3) и
фазы поля $\Phi$. После калибровочного факторирования $SO(3)$ и фазового
U(1) остаются семь строго положительных собственных значений. Минимальное
из них равно
\begin{equation}
 \frac{32-8\sqrt7}{3}=3.61132984\ldots
\end{equation}
Отрицательных направлений нет.

\section{Единичный дефектный импульс}

Добавим уже выведенный на корневой ветви вклад
\begin{equation}
 V_k=|\Phi|^2.
\end{equation}
Точная конденсированная стационарная точка имеет вид
\begin{equation}
 X_\star=\sqrt{\frac56}\,I_3,
 \qquad
 |\Phi_\star|=\sqrt{\frac23},
 \qquad
 V_\star=\frac56.
\end{equation}
Нормальная ветвь имеет энергию $3/2$, поэтому конденсированная ветвь ниже
на
\begin{equation}
 \Delta V=\frac23.
\end{equation}
Полный одиннадцатимерный спектр второй вариации равен
\begin{equation}
 \left\{
 0^{(4)},
 \left(\frac{40}{9}\right)^{(5)},
 \frac{68-4\sqrt{109}}9,
 \frac{68+4\sqrt{109}}9
 \right\}.
 \label{eq:v5-m300-momentum-hessian}
\end{equation}
После того же калибровочного факторирования физический скалярный гессиан
строго положителен. Масса каждой из двух связующих комбинаций удовлетворяет
\begin{equation}
 m_{\mathrm{conn}}^2=\frac32.
\end{equation}

\begin{theorem}[Устойчивость связующего сектора]
Единоследовое действие~\eqref{eq:v5-m300-hodge-action} имеет устойчивый
конденсированный минимум как при нулевом, так и при единичном дефектном
импульсе. После удаления калибровочных нулевых направлений все скалярные
собственные значения положительны.
\end{theorem}

\section{Граница единственности действия}

Полный неразмеченный след $M_{300}$ фиксирует равные веса трёх блоков
кривизны в~\eqref{eq:v5-m300-hodge-action}. Однако меньшая блочная симметрия
сама по себе допускает по меньшей мере отдельные веса крайнего и среднего
блоков. После удаления общего масштаба остаётся один относительный
коэффициент.

Следовательно, действие не является следствием одной блочной симметрии.
Оно является отдельным родительским принципом: используется одна
невзвешенная норма полной кривизны. Этот принцип зафиксирован до
феноменологии и не подбирался по массам, но его необходимо проверить на
полном пространстве допустимых одноформ.

\section{Почему полный гессиан ещё не закрыт}

Вычисленный спектр охватывает все компоненты $(X,\Phi)$ и их калибровочные
орбиты, но ещё не охватывает:
\begin{enumerate}
 \item возможные одноформы между 15-мерным равномерным опорным сектором и
 цепью;
 \item все разрешённые смешивания углов $M_{60}$ и $M_{270}$;
 \item поля БВ/БРСТ, духи и якобианы полного представления;
 \item семейный проектор ранга один и юкавские флуктуации в том же
 функционале.
\end{enumerate}
Само действие~\eqref{eq:v5-m300-hodge-action} не поднимает опорный сектор,
если тот объявлен отдельным нулевым блоком. Его удаление пока опирается на
канонический проектор $P_3$, а не на полный физический фактор.

\begin{nogocriterion}[Неполный гессиан $M_{300}$]
Положительность одиннадцатимерного связующего гессиана не считается полной
физической устойчивостью, пока не классифицировано всё пространство
представленных одноформ $M_{300}$ и не доказано, что опорные и смешанные
направления либо отсутствуют по условию первого порядка, либо получают
положительную массу из того же квадрата без второго веса.
\end{nogocriterion}

\section{Вердикт}

\begin{protocol}
\begin{enumerate}
 \item один квадрат кривизны и один след:
 \textbf{получены};
 \item независимые портальные коэффициенты внутри кандидата:
 \textbf{отсутствуют};
 \item точный связующий гессиан при нулевом импульсе:
 \textbf{положителен после факторирования};
 \item устойчивый конденсат при единичном импульсе:
 \textbf{получен};
 \item масса двух связующих комбинаций:
 \textbf{положительна};
 \item единственность действия из одной блочной симметрии:
 \textbf{не доказана};
 \item полное пространство одноформ $M_{300}$:
 \textbf{не классифицировано};
 \item полный физический гессиан БВ/БРСТ:
 \textbf{не получен};
 \item физическое замыкание:
 \textbf{не достигнуто}.
\end{enumerate}
\end{protocol}

Следующая проверка должна классифицировать полный комплекс представленных
одноформ и смешанные опорно-цепные блоки $M_{300}$, после чего повторить
гессиан с БВ/БРСТ-факторированием.

Машинная проверка и сертификат сохранены в
\path{s2t_v5_m300_hodge_curvature_hessian_gate.py} и
\path{s2t_v5_m300_hodge_curvature_hessian_gate_results.json}.